\documentclass[12pt,a4paper]{article}
\usepackage[margin=2.5cm]{geometry}
\usepackage{booktabs}
\usepackage{longtable}
\usepackage{array}
\usepackage{xcolor}
\usepackage{hyperref}
\usepackage{listings}
\usepackage{amsmath}

\hypersetup{colorlinks=true, linkcolor=blue}

\lstdefinestyle{pseudo}{
  basicstyle=\ttfamily\small,
  breaklines=true,
  frame=single,
  backgroundcolor=\color{gray!10},
  columns=flexible,
  mathescape=true,
}

\title{\textbf{SLR(1) Parse Table}\\
       \large Decaf Mini-Compiler\\
       CS-471L Compiler Construction Lab --- UET Lahore, Spring 2026}
\author{Compiler Project Team}
\date{}

\begin{document}
\maketitle
\tableofcontents
\newpage

% ------------------------------------------------------------------
\section{Algorithm Choice: SLR(1)}
% ------------------------------------------------------------------

We chose SLR(1) for the following reasons:

\begin{enumerate}
  \item \textbf{More powerful than LR(0)}: FOLLOW sets restrict reduce actions,
        eliminating spurious conflicts.
  \item \textbf{Simpler than LALR(1)/LR(1)}: no lookahead propagation is needed
        through item sets.
  \item \textbf{Compatible with the Decaf grammar}: with four explicit conflict
        resolutions, the grammar is SLR(1)-parseable.
  \item \textbf{Reuses FOLLOW sets}: already computed for the LL(1) analysis.
\end{enumerate}

The Decaf SLR(1) automaton has \textbf{229 states} and \textbf{4 resolved
conflicts}.

% ------------------------------------------------------------------
\section{Bottom-Up Parsing Concepts}
% ------------------------------------------------------------------

\subsection{LR(0) Items}

An LR(0) item is a production with a dot marker:
\[
  [A \to \alpha \bullet \beta]
\]
The dot shows how much of the right-hand side has been matched.
The dot advances on each shift or goto action.

\subsection{Closure}

\begin{lstlisting}[style=pseudo]
closure(I):
  result = I
  for each [A -> $\alpha$ . B $\beta$] in result:
    for each production B -> $\gamma$:
      add [B -> . $\gamma$] to result
  return result
\end{lstlisting}

\subsection{Goto}

\begin{lstlisting}[style=pseudo]
goto(I, X):
  moved = { [A -> $\alpha$ X . $\beta$] : [A -> $\alpha$ . X $\beta$] in I }
  return closure(moved)
\end{lstlisting}

\subsection{SLR(1) Table Construction}

\begin{lstlisting}[style=pseudo]
for each state $s$ and each item in state $s$:

  [A -> $\alpha$ . a $\beta$]  (a is a terminal):
    action[$s$, $a$] = shift goto($s$, $a$)

  [A -> $\alpha$ . B $\beta$]  (B is a non-terminal):
    goto[$s$, $B$] = goto($s$, $B$)

  [A -> $\alpha$ .]:
    for each $t \in$ FOLLOW($A$):
      action[$s$, $t$] = reduce ($A \to \alpha$)

  [S' -> S .]:
    action[$s$, $\$$] = accept
\end{lstlisting}

% ------------------------------------------------------------------
\section{Key LR(0) Item Sets}
% ------------------------------------------------------------------

\subsection{State 0: Initial State}

\begin{lstlisting}[style=pseudo]
[Program' -> . Program]
[Program  -> . DeclList]
[DeclList -> . Decl DeclList]
[DeclList -> .]                      -- reduce on $
[Decl     -> . VarDecl]
[Decl     -> . FuncDecl]
[Decl     -> . ClassDecl]
[Decl     -> . InterfaceDecl]
[VarDecl  -> . Type id ;]
[FuncDecl -> . RetType id ( Formals ) Block]
[Type     -> . int TypeSuffix]
[Type     -> . double TypeSuffix]
[Type     -> . bool TypeSuffix]
[Type     -> . string TypeSuffix]
[Type     -> . id TypeSuffix]
[RetType  -> . Type]
[RetType  -> . void]
[ClassDecl -> . class id ...]
[InterfaceDecl -> . interface id ...]
\end{lstlisting}

\subsection{State 108: TypeSuffix Conflict}

\begin{lstlisting}[style=pseudo]
[TypeSuffix -> [ ] . TypeSuffix]    -- can shift [ for nested suffix
[TypeSuffix -> .]                   -- can reduce TypeSuffix -> epsilon
\end{lstlisting}

Conflict on \texttt{[}: shift vs reduce. Resolution: \textbf{reduce} (type
context is complete).

\subsection{State 218: Dangling Else Conflict}

\begin{lstlisting}[style=pseudo]
[IfStmt -> if ( Expr ) Stmt . ElsePart]
[ElsePart -> .]                          -- reduce on FOLLOW(ElsePart)
[ElsePart -> . else Stmt]               -- shift on else
\end{lstlisting}

Conflict on \texttt{else}: shift vs reduce. Resolution: \textbf{shift}
(standard dangling-else binding).

% ------------------------------------------------------------------
\section{Action Table: Selected Entries}
% ------------------------------------------------------------------

Notation: $s_N$ = shift to state $N$; $r_P$ = reduce by production $P$;
acc = accept.

\begin{center}
\begin{tabular}{l|ccccccc}
\toprule
State & \texttt{int} & \texttt{void} & \texttt{class} & \texttt{id} & \texttt{\{} & \texttt{\$} & \texttt{else} \\
\midrule
0   & $s_{10}$ & $s_{15}$ & $s_{16}$ & $s_{14}$ & --- & --- & --- \\
1   & ---      & ---      & ---      & $r$(resolve) & --- & acc & --- \\
2   & ---      & ---      & ---      & ---      & --- & $r_0$ & --- \\
10  & ---      & ---      & ---      & $r_{17}$ & --- & $r_{17}$ & --- \\
218 & $r_{55}$ & $r_{55}$ & $r_{55}$ & $r_{55}$ & $r_{55}$ & $r_{55}$ & $s_{219}$ \\
\bottomrule
\end{tabular}
\end{center}

\textit{State 218, \texttt{else}: shift to state 219 (resolve dangling else by shifting).}

% ------------------------------------------------------------------
\section{Conflict Resolution Summary}
% ------------------------------------------------------------------

\begin{center}
\begin{tabular}{llllp{5cm}}
\toprule
\# & State & Token & Type & Resolution \\
\midrule
1 & 218 & \texttt{else} & Shift/Reduce & Prefer shift; binds else to nearest if \\
2 & 108 & \texttt{[}    & Shift/Reduce & Prefer reduce; type context complete \\
3 & 108 & \texttt{)}    & Reduce/Reduce & Use first matching production \\
4 & 1   & \texttt{id}   & Shift/Reduce & Prefer reduce; complete DeclList \\
\bottomrule
\end{tabular}
\end{center}

% ------------------------------------------------------------------
\section{SLR(1) Parsing Algorithm}
% ------------------------------------------------------------------

\begin{lstlisting}[style=pseudo]
state_stack = [0]
symbol_stack = [$]

while True:
  s = state_stack[-1]
  a = next_token.grammar_symbol()

  if action[s][a] == shift t:
    symbol_stack.push(a)
    state_stack.push(t)
    advance_input()

  elif action[s][a] == reduce (A -> beta):
    for _ in range(len(beta)):
      symbol_stack.pop()
      state_stack.pop()
    t = goto[state_stack[-1]][A]
    symbol_stack.push(A)
    state_stack.push(t)

  elif action[s][a] == accept:
    return SUCCESS

  else:
    error_handler.report(s, a)
    panic_mode_recover(FOLLOW(current_lhs))
\end{lstlisting}

% ------------------------------------------------------------------
\section{Parse Trace: test1\_valid.decaf (Excerpt)}
% ------------------------------------------------------------------

Input: \texttt{void main() \{ int x; x = 5; Print(z); \}}

\begin{longtable}{rp{3cm}p{3.5cm}p{3cm}p{3cm}}
\toprule
\textbf{Step} & \textbf{State Stack} & \textbf{Symbol Stack} & \textbf{Input} & \textbf{Action} \\
\midrule
\endhead
\bottomrule
\endfoot
1  & [0] & [\$] & void main... & shift void \\
2  & [0,15] & [\$,void] & main... & shift id \\
3  & [0,15,N] & [\$,void,id] & ( )... & goto RetType \\
4  & [0,...,R] & [\$,RetType] & main... & shift id=main \\
5  & [...] & [\$,RetType,id] & (... & shift ( \\
6  & [...] & [\$,RetType,id,(] & )... & $r_{19}$ Formals$\to\varepsilon$ \\
7  & [...] & [\$,RetType,id,(,Formals] & )... & shift ) \\
8  & [...] & [\$,...,)] & \{... & goto Block \\
9  & [...] & [\$,...,Block...] & \{... & shift \{ \\
10 & [...] & [\$,...,\{] & int x;... & shift int \\
11 & [...] & [\$,...,int] & x;... & $r_{17}$ TypeSuffix$\to\varepsilon$ \\
12 & [...] & [\$,...,TypeSuffix] & x;... & $r_{11}$ Type$\to$int TypeSuffix \\
13 & [...] & [\$,...,Type] & x;... & shift id=x \\
14 & [...] & [\$,...,Type,id] & ;... & shift ; \\
15 & [...] & [\$,...,Type,id,;] & x=5... & $r_7$ VarDecl$\to$Type id ; \\
16 & [...] & [\$,...,VarDecl] & x=5... & $r_{43}$ Stmt$\to$VarDecl \\
17 & [...] & [\$,...,Stmt] & x=5... & goto StmtList \\
18 & [...] & [\$,...,StmtList] & x=5... & shift id=x (ExprStmt) \\
19 & [...] & [\$,...,id] & =5... & shift = \\
20 & [...] & [\$,...,=] & 5;... & ... (reduce expr chain) \\
\end{longtable}

Full trace available in \texttt{output/test1\_valid\_lr\_parser.txt}.

\end{document}
